% ==============================================================================
% Related Work and Methodology Draft for Confidence-Gated Code Verification
% Formatted for IEEE / ACM Double-Column Format
% ==============================================================================

\section{Related Work}
\label{sec:related_work}

Our work bridges four emerging frontiers in machine learning, uncertainty quantification, and software verification:

\subsection{Lightweight and Resource-Constrained Hallucination Detection}
Hallucination detection in large language models (LLMs) has predominantly relied on heavy computational paradigms, such as multi-agent consensus debate, fine-tuned 70B parameter LLM-as-a-judge evaluators, or high-dimensional hidden-state probing. Recently, Faujdar and Kadvani~\cite{faujdar2026gpu} investigated the practical boundaries of resource-constrained deployment in \emph{``How Far Can You Get Without a GPU?''}, demonstrating that lightweight, CPU-feasible feature ensembles (e.g., token overlap, lightweight cross-encoders, and NLI classifiers) can match GPU-bound detectors on closed-domain benchmarks. 

Concurrently, Nguyen et al.~\cite{nguyen2026probabilities} challenged the necessity of multi-sample generation in \emph{``Probabilities Are All You Need: A Probability-Only Approach to Uncertainty Estimation''}, proving that single-pass output probabilities over top-$K$ candidate tokens can approximate predictive entropy without incurring the latency and compute of sampling multiple responses ($N=10$). 

Our work adapts this lightweight, single-pass philosophy to program synthesis. Unlike free-form natural language, code exhibits formal Context-Free Grammar (CFG) structures and deterministic call interfaces. We exploit this property to extract decision-point token entropy and Abstract Syntax Tree (AST) structural features in a single forward pass, providing zero-GPU verification within milliseconds on consumer hardware.

\subsection{Uncertainty Estimation and Beyond-Entropy Signals}
Uncertainty quantification via Shannon entropy and semantic equivalence clustering~\cite{kuhn2023semantic} has emerged as the standard approach for detecting factual hallucinations. However, post-softmax probability distributions frequently suffer from poor calibration and overconfidence. 

In \emph{``Semantic Energy: Detecting LLM Hallucination Beyond Entropy''}, Ma et al.~\cite{ma2025semantic} demonstrated that standard semantic entropy collapses when an LLM operates in an overconfident regime, proposing unnormalized energy metrics from intermediate representations to capture epistemic uncertainty. Similarly, Lee and Yoshinaga~\cite{lee2026logit} showed that analyzing layer-wise semantic logit trajectories via a compact Multi-Layer Perceptron (MLP) effectively discriminates hallucinations without storing or processing dense 4096-dimensional hidden states.

In this work, we uncover a critical manifestation of this phenomenon in code generation: the \emph{``Confidently Wrong''} paradox. While token-level Shannon entropy effectively separates general syntactic degradation ($\text{AUROC} \approx 0.855$), it completely inverts when applied to usage-semantic intent misuse ($\text{AUROC} \approx 0.389$). Because semantic hallucinations stem from false interface beliefs rather than decoding hesitation, models generate contract-violating code with near-zero token entropy.

\subsection{Adaptive Token Selection and Localized Probing}
A major limitation of sequence-level uncertainty metrics is token dilution: averaging predictive entropy across hundreds of boilerplate tokens (indentation, keywords, trivial operators) washes out the critical error signal. 

To mitigate this in unstructured NLP, Niu, Haddadi, and Pang~\cite{niu2025robust} introduced \emph{HaMI} (\emph{``Robust Hallucination Detection in LLMs via Adaptive Token Selection''}, NeurIPS 2025), formulating hallucination detection as Multiple Instance Learning (MIL) over token representations to isolate the sparse subset of tokens driving factual divergence. In parallel, Sun et al.~\cite{sun2026adaptive} (AAAI 2026) developed an adaptive Bayesian estimation framework to dynamically adjust sampling budgets based on variance convergence.

In program synthesis, token selection does not require heuristic statistical pooling. Because source code conforms to formal grammars, we introduce deterministic \emph{Decision-Point Alignment}: token character offsets are aligned directly with AST Call and Def-Use sites, isolating predictive entropy precisely at the interface where API return values are manipulated.

\subsection{Formal Verification and Theoretical Impossibility Limits}
Can purely statistical models ever achieve sound hallucination detection? In \emph{``Is Automated Hallucination Detection Fundamentally Possible?''}, Karbasi et al.~\cite{karbasi2025impossibility} (Yale University) established a theoretical equivalence between hallucination detection and classical language identification in the limit~\cite{gold1967language,angluin1980inductive}. They proved that black-box statistical learners cannot guarantee bounded detection error over unconstrained domains without access to formal structural verification oracles.

This theoretical boundary directly explains why pure neural or entropy scoring fails on semantic code misuse. Our architecture introduces a \emph{Neuro-Symbolic} cascade: statistical machine learning models (Random Forests, Gradient Boosted Decision Trees, and Focal-Loss Neural Networks) filter the probabilistic fluency of code, while formal AST Def-Use contract checkers enforce deterministic invariant boundaries over known library interfaces.

% ==============================================================================
% SECTION: METHODOLOGY
% ==============================================================================

\section{Methodology}
\label{sec:methodology}

\subsection{Problem Formulation and Failure Taxonomy}
Let $x$ denote an input programming specification (docstring and function signature), and let $\mathbf{y} = (y_1, y_2, \dots, y_T)$ denote an autoregressive code sequence generated by an LLM $P_\theta(y_t \mid y_{<t}, x)$. At each generation step $t$, the model emits top-$K$ candidate tokens and unnormalized log-probabilities $(w_{t,k}, p_{t,k})_{k=1}^K$. The normalized Shannon entropy at step $t$ is computed as:
\begin{equation}
    H(t) = - \sum_{k=1}^K \tilde{p}_{t,k} \log_2 \tilde{p}_{t,k}, \quad \text{where } \tilde{p}_{t,k} = \frac{p_{t,k}}{\sum_{j=1}^K p_{t,j}}
\end{equation}

We define two ground-truth evaluation targets based on sandboxed unit test execution:
\begin{enumerate}
    \item \textbf{General Failure ($Y_{\text{fail}} \in \{0, 1\}$)}: Execution failure encompassing syntax errors, uncaught runtime exceptions, and assertion mismatches.
    \item \textbf{Usage-Semantic Intent Misuse ($Y_{\text{misuse}} \in \{0, 1\}$)}: Grammatically sound code that invokes valid third-party API identifiers but violates the dynamic return-type contract or lifecycle assumptions of the library.
\end{enumerate}

\subsection{The Confidence-Gated Verification Cascade}
The verification pipeline comprises three sequential layers designed to minimize verification latency and compute:

\subsubsection{Layer 1: Near-Zero-Cost Feature Extraction}
From completion $\mathbf{y}$, we extract a 21-dimensional feature vector $\mathbf{z} \in \mathbb{R}^{21}$ across four domains:
\begin{itemize}
    \item \textbf{Decision-Point Entropy (6 features)}: Mean $\bar{H}_{dp}$, maximum $\max H_{dp}$, and standard deviation $\sigma(H_{dp})$ over tokens immediately succeeding AST Call and Def-Use sites; entropy contrast $\Delta H = \bar{H}_{dp} - \bar{H}_{seq}$; entropy ratio $R_H$; and decision token fraction $r_{tok}$.
    \item \textbf{Sequence Uncertainty (4 features)}: Sequence-wide mean $\bar{H}_{seq}$, maximum $\max H_{seq}$, standard deviation $\sigma(H_{seq})$, and high-entropy token fraction $\%H_{>2.0}$.
    \item \textbf{Cross-Sample Usage Diversity (1 feature)}: Shannon diversity $H_{div}$ over downstream API consumption patterns across $N$ generations.
    \item \textbf{AST Complexity \& Contracts (10 features)}: AST depth, total node count, code character length, token count, API call count; and 5 explicit AST Def-Use indicators ($I_{misuse}$, $C_{misuse}$, $N_{def\_use}$, $I_{subscript}$, $I_{type\_invalid}$).
\end{itemize}

\subsubsection{Layer 2: Fast Confidence Router}
A lightweight ensemble classifier (Random Forest or HistGradientBoosting) evaluates $\mathbf{z}$ to produce a calibrated failure probability $p_{\text{cheap}} \in [0, 1]$. Triage decisions follow dual confidence thresholds:
\begin{equation}
    \text{Triage}(p_{\text{cheap}}) = \begin{cases}
        \text{FAST\_ACCEPT} & \text{if } p_{\text{cheap}} \le \tau_{\text{accept}} \\
        \text{FAST\_REJECT} & \text{if } p_{\text{cheap}} \ge \tau_{\text{reject}} \\
        \text{ESCALATE}     & \text{if } \tau_{\text{accept}} < p_{\text{cheap}} < \tau_{\text{reject}}
    \end{cases}
\end{equation}
By setting $\tau_{\text{accept}} = 0.30$ and $\tau_{\text{reject}} = 0.75$, the router settles 93\% of completions with near-zero latency, circumventing test sandboxing for confident passes and obvious syntax collapse.

\subsubsection{Layer 3: Distilled Specialist Verifier}
Borderline cases $(\tau_{\text{accept}} < p_{\text{cheap}} < \tau_{\text{reject}})$ are escalated to Layer 3, an MLP verifier trained with Binary Focal Loss ($\gamma = 2.0$) on multi-modal inputs combining tabular signals $\mathbf{z}$ and sub-token TF-IDF $N$-grams $\mathbf{c} \in \mathbb{R}^{48}$:
\begin{equation}
    \mathcal{L}_{\text{FL}}(p_t) = -\alpha_t (1 - p_t)^\gamma \log(p_t)
\end{equation}
The focusing parameter $\gamma = 2.0$ suppresses easy negative gradients, forcing the specialist head to resolve subtle semantic contract violations.

% ==============================================================================
% SECTION: EXPERIMENTAL RESULTS
% ==============================================================================

\section{Experimental Evaluation}
\label{sec:experiments}

\begin{table*}[t]
\centering
\small
\caption{Out-of-fold generalization performance across classifier architectures under 5-fold GroupKFold cross-validation grouped by task identifier (zero task leakage). Machine learning models without rules capture general failure effectively ($\text{AUROC} \approx 0.865$), while AST Def-Use contract indicators resolve the semantic misuse blind spot.}
\label{tab:out_of_fold_benchmark}
\begin{tabular}{l|ccc|ccc}
\hline
& \multicolumn{3}{c|}{\textbf{General Failure Detection ($Y_{\text{fail}}$)}} & \multicolumn{3}{c}{\textbf{Intent Misuse Detection ($Y_{\text{misuse}}$)}} \\
\textbf{Model Architecture} & \textbf{AUROC} & \textbf{AUPRC} & \textbf{F1 Score} & \textbf{AUROC} & \textbf{AUPRC} & \textbf{F1 Score} \\
\hline
Logistic Regression (Balanced)   & $0.778 \pm 0.03$ & $0.906 \pm 0.08$ & $0.736$ & $\mathbf{1.000 \pm 0.00}$ & $\mathbf{1.000 \pm 0.00}$ & $\mathbf{1.000}$ \\
Random Forest (100 Trees)        & $0.757 \pm 0.05$ & $0.892 \pm 0.08$ & $0.825$ & $\mathbf{1.000 \pm 0.00}$ & $\mathbf{1.000 \pm 0.00}$ & $0.778$ \\
HistGradientBoosting             & $0.855 \pm 0.06$ & $0.952 \pm 0.03$ & $\mathbf{0.845}$ & $0.556 \pm 0.33$ & $0.431 \pm 0.40$ & $0.000$ \\
Multi-Layer Perceptron (MLP)     & $\mathbf{0.865 \pm 0.05}$ & $\mathbf{0.960 \pm 0.02}$ & $0.773$ & $0.792 \pm 0.30$ & $0.731 \pm 0.38$ & $0.667$ \\
\hline
Baseline (Naive Prevalence)      & — & ($0.750$) & — & — & ($0.117$) & — \\
\hline
\end{tabular}
\end{table*}

\begin{table}[t]
\centering
\small
\caption{Confidence Cascade Tradeoff Sweep on Intent Misuse ($Y_{\text{misuse}}$) under leak-free out-of-fold evaluation. The balanced cascade reduces verification compute by 93\% while maintaining optimal out-of-fold classification accuracy.}
\label{tab:cascade_sweep}
\begin{tabular}{l|cccc}
\hline
\textbf{Routing Strategy} & \textbf{Escalation \%} & \textbf{Cost Saved \%} & \textbf{Accuracy \%} & \textbf{F1 Score} \\
\hline
Pure Cheap Router (0\% Escalate)     & 0.0\%   & 100.0\% & 97.0\%  & 0.727 \\
Always-Escalate (100\% Specialist)   & 100.0\% & 0.0\%   & 100.0\% & 1.000 \\
Confidence Cascade (Balanced)        & 7.0\%   & 93.0\%  & 100.0\% & 1.000 \\
Confidence Cascade (High-Efficiency) & 5.0\%   & 95.0\%  & 98.0\%  & 0.833 \\
\hline
\end{tabular}
\end{table}

\subsection{Empirical Analysis & Statistical Framing}
As documented in Table~\ref{tab:out_of_fold_benchmark} and Table~\ref{tab:cascade_sweep}, our balanced cascade reduces sandbox execution requirements by 93\% while achieving an empirical F1-score of 1.000 on Intent Misuse. 

We explicitly contextualize this metric within its statistical boundaries:
In our 100-sample pilot benchmark, ground-truth Intent Misuse constitutes 7\% of samples ($N=7$). Our AST Def-Use semantic rules successfully flag all 7 samples with zero false positives on the 93 non-misuse completions, yielding an empirical recall of 100\% with a 95\% Wilson score confidence interval of $[64.6\%, 100.0\%]$. Rather than asserting a generalized 100\% guarantee over all future code, this demonstrates that \emph{anchoring probabilistic token signals with deterministic AST def-use tracking resolves the empirical blind spot of semantic entropy}, providing a scalable verification framework for production software systems.
